% S8. Qué evidencia hay y qué no hay de que Binance sufra este bucle.
\begin{tikzpicture}[x=1cm,y=1cm]
  \path[use as bounding box] (0,0.1) rectangle (13.5,4.6);
  \tikzset{item/.style={box,font=\scriptsize,text width=5.7cm,align=left,anchor=west}}
  \node[anchor=west,font=\small\bfseries] at (0.1,4.35) {Lo que sí hay: indicios};
  \node[item] at (0.1,3.45) {Usuarios del \textbf{foro oficial} de Binance reportan conexiones cerradas con el código \textbf{1006} durante trading intenso.};
  \node[item] at (0.1,2.3) {También reportan \textbf{muchos streams de profundidad} cerrándose al mismo tiempo.};
  \node[item] at (0.1,1.15) {\textbf{Tardis.dev}, proveedor de datos, reconstruye el libro con el \textbf{snapshot REST} cuando pierde un mensaje.};
  \draw[colNeutro!40] (6.75,0.15) -- (6.75,4.55);
  \node[anchor=west,font=\small\bfseries,text=colWake] at (7.1,4.35) {Lo que no hay};
  \node[item,malo,align=left] at (7.1,3.45) {Ningún \textbf{post-mortem público} de Binance atribuye una caída a este bucle.};
  \node[item,malo,align=left] at (7.1,2.3) {No conocemos sus \textbf{mitigaciones internas} ni su capacidad real.};
  \node[item,malo,align=left] at (7.1,1.15) {No sabemos \textbf{cuánto le cuesta generar un snapshot}: lo variamos en la simulación.};
\end{tikzpicture}
